%; whizzy paragraph -pdf xpdf -latex ./whizzypdfptex.sh
%; whizzy-paragraph "^\\\\begin{frame}\\|\\\\emtext"
% latex beamer presentation.
% platex, latex-beamer でコンパイルすることを想定。 

%     Tokyo Debian Meeting resources
%     Copyright (C) 2012 Junichi Uekawa

%     This program is free software; you can redistribute it and/or modify
%     it under the terms of the GNU General Public License as published by
%     the Free Software Foundation; either version 2 of the License, or
%     (at your option) any later version.

%     This program is distributed in the hope that it will be useful,
%     but WITHOUT ANY WARRANTY; without even the implied warreanty of
%     MERCHANTABILITY or FITNESS FOR A PARTICULAR PURPOSE.  See the
%     GNU General Public License for more details.

%     You should have received a copy of the GNU General Public License
%     along with this program; if not, write to the Free Software
%     Foundation, Inc., 51 Franklin St, Fifth Floor, Boston, MA  02110-1301 USA

\documentclass[cjk,dvipdfmx,12pt]{beamer}
\usetheme{Tokyo}
\usepackage{monthlypresentation}

%  preview (shell-command (concat "evince " (replace-regexp-in-string "tex$" "pdf"(buffer-file-name)) "&")) 
%  presentation (shell-command (concat "xpdf -fullscreen " (replace-regexp-in-string "tex$" "pdf"(buffer-file-name)) "&"))
%  presentation (shell-command (concat "evince " (replace-regexp-in-string "tex$" "pdf"(buffer-file-name)) "&"))

%http://www.naney.org/diki/dk/hyperref.html
%日本語EUC系環境の時
\AtBeginDvi{\special{pdf:tounicode EUC-UCS2}}
%シフトJIS系環境の時
%\AtBeginDvi{\special{pdf:tounicode 90ms-RKSJ-UCS2}}

\newenvironment{commandlinesmall}%
{\VerbatimEnvironment
  \begin{Sbox}\begin{minipage}{1.0\hsize}\begin{fontsize}{8}{8} \begin{BVerbatim}}%
{\end{BVerbatim}\end{fontsize}\end{minipage}\end{Sbox}
  \setlength{\fboxsep}{8pt}
% start on a new paragraph

\vspace{6pt}% skip before
\fcolorbox{dancerdarkblue}{dancerlightblue}{\TheSbox}

\vspace{6pt}% skip after
}
%end of commandlinesmall

\title{東京エリアDebian勉強会}
\subtitle{第133回 2015年11月度}
\author{野島貴英}
\date{2015年11月21日}
\logo{\includegraphics[width=8cm]{image200607/openlogo-light.eps}}

\begin{document}

\begin{frame}
\titlepage{}
\end{frame}

\begin{frame}{Agenda}
 \begin{minipage}[t]{0.45\hsize}
  \begin{itemize}
   \item 事前課題発表
   \item 最近あったDebian関連のイベント報告
	 \begin{itemize}
	 \item 第131回東京エリアDebian勉強会
         \item 第132回東京エリアDebian勉強会(OSC 2015 Tokyo/Fall)
	 \end{itemize}
  \end{itemize}
 \end{minipage} 
 \begin{minipage}[t]{0.45\hsize}
  \begin{itemize}
   \item Debian Trivia Quiz
   \item Debian GNU/kFreeBSD セットアップガイド 2015年版
   \item 今後のイベント
   \item 今日の宴会場所
  \end{itemize}
 \end{minipage}
\end{frame}

\section{事前課題}
\emtext{事前課題}
{\footnotesize
 \begin{prework}{ wskoka }
  \begin{enumerate}
  \item Q.hack time $B$K2?$r$7$^$9$+!)(B\\
    A. Tilera$B$N(BDebian$B:n$j$?$$!#(B
  \item Q.$BK\JY6/2q$r$I$3$G$*CN$j$K$J$j$^$7$?$+!)(B\\
    A. dots$B$N%5%$%HWQWKCf$KH/8+!#(B
  \end{enumerate}
\end{prework}

\begin{prework}{ henrich }
  \begin{enumerate}
  \item Q.hack time $B$K2?$r$7$^$9$+!)(B\\
    A. $BK]Lu:n6H$r?J$a$?$$$H;W$$$^$9!#(B
  \end{enumerate}
\end{prework}

\begin{prework}{ hiroyuki.nagata }
  \begin{enumerate}
  \item Q.hack time $B$K2?$r$7$^$9$+!)(B\\
    A. $B$J$K$+%W%m%0%i%`$r=q$/(B
  \item Q.$BK\JY6/2q$r$I$3$G$*CN$j$K$J$j$^$7$?$+!)(B\\
    A. $B$?$7$+(BTwitter
  \end{enumerate}
\end{prework}

\begin{prework}{ dictoss }
  \begin{enumerate}
  \item Q.hack time $B$K2?$r$7$^$9$+!)(B\\
    A. Debian GNU/kFreeBSD$B$NF0:n%F%9%H(B
  \item Q.$BK\JY6/2q$r$I$3$G$*CN$j$K$J$j$^$7$?$+!)(B\\
    A. $B8}%3%_(B
  \end{enumerate}
\end{prework}

\begin{prework}{ rosh }
  \begin{enumerate}
  \item Q.hack time $B$K2?$r$7$^$9$+!)(B\\
    A. BTS$B$N:n6H(B
  \end{enumerate}
\end{prework}

\begin{prework}{ koedoyoshida }
  \begin{enumerate}
  \item Q.hack time $B$K2?$r$7$^$9$+!)(B\\
    A. $B$"$s$I$-$e$a$s$F$C$I$G$S$"$sE_9f$NJT=8(B
  \item Q.$BK\JY6/2q$r$I$3$G$*CN$j$K$J$j$^$7$?$+!)(B\\
    A. $BA02s$NJY6/2q(B
  \end{enumerate}
\end{prework}

\begin{prework}{ yy\_y\_ja\_jp }
  \begin{enumerate}
  \item Q.hack time $B$K2?$r$7$^$9$+!)(B\\
    A. DDTSS
  \end{enumerate}
\end{prework}

\begin{prework}{ $B$+$8(B }
  \begin{enumerate}
  \item Q.hack time $B$K2?$r$7$^$9$+!)(B\\
    A. XCB$B%W%m%0%i%_%s%0(B
  \end{enumerate}
\end{prework}

\begin{prework}{ $BLnEg(B }
  \begin{enumerate}
  \item Q.hack time $B$K2?$r$7$^$9$+!)(B\\
    A. DDTSS$B$H$+(B
  \item Q.$BK\JY6/2q$r$I$3$G$*CN$j$K$J$j$^$7$?$+!)(B\\
    A. $B44;v$@$C$?$j$G$7$F(B...
  \end{enumerate}
\end{prework}




}

\section{イベント報告}
\emtext{イベント報告}

\begin{frame}{第131回東京エリアDebian勉強会 }

\begin{itemize}
\item 場所はdotsさんをお借りしての開催でした。
\item 参加者は7名でした。
\item セミナ内容は野島さんによる「DebConf15ビデオ紹介」、Rogerさんによる「毎日使えるIPv6ネットワークの構築」でした。
\item 残りの時間でhack timeを行い、成果発表をしました。
\end{itemize} 
\end{frame}

\begin{frame}{第131回東京エリアDebian勉強会(つづき)}

　野島さんによる「DebConf15ビデオ紹介」は、DebConf15の公開ビデオから、
\begin{itemize}
  \item ``Stretching out for trustworthy reproducible builds''
  \item ``Thanks for maintaining a desktop environment. But is it accessible?''
  \item ``Citizenfour Screening''
\end{itemize}
についての、サマリとハイライトについて語りました。

\end{frame}

\begin{frame}{第131回東京エリアDebian勉強会(つづき)}

 　Rogerさんによる「毎日使えるIPv6ネットワークの構築」は、実際に自宅にDebian機をIPv6環境でつなぐことに関する苦労話や、設定方法についての発表でした。昨今、日本の総務省からも、各携帯3社へ、IPv6接続の関する圧力がかかる状況\footnote{http://www.itmedia.co.jp/news/articles/1511/11/news062.html}ですので、Debianの熱烈ユーザの皆様におきましては、是非こちらの発表をきっかけにIPv6接続をDebianでお試しください。
  
\end{frame}

\begin{frame}{第132回東京エリアDebian勉強会(OSC 2015 Tokyo/Fall)}

 \begin{itemize}
\item 場所は明星大学さんでOSC 2015 Tokyo/Fallのイベントが行われました。
\item 展示用のブースを出し、Debianをこれから初める方、Debianを使って困っている方などいろいろな方から意見、議論をして交流を行いました。
\item セミナの内容は岩松さんによる「Debian と systemd について」でした。
\item ブースには65名、セミナは部屋がほぼ満員でした。
\end{itemize} 
  
\end{frame}

\section{Debian Trivia Quiz}
\emtext{Debian Trivia Quiz}
\begin{frame}{Debian Trivia Quiz}

  Debian の常識、もちろん知ってますよね?
知らないなんて恥ずかしくて、知らないとは言えないあんなことやこんなこと、
みんなで確認してみましょう。

今回の出題範囲は\url{debian-devel-announce@lists.debian.org},
\url{debian-news@lists.debian.org} に投稿された
内容などからです。

\end{frame}

\subsection{問題}

%; whizzy-master ../debianmeetingresume201311.tex
% $B0J>e$N@_Dj$r$7$F$$$k$?$a!"$3$N%U%!%$%k$G(B M-x whizzytex $B$9$k$H!"(Bwhizzytex$B$,MxMQ$G$-$^$9!#(B
%

\santaku
{2015/10/22$B$N(BDPN$B$N%a!<%k$+$iDj4|E*$K$J$,$l$F$$$?$$$/$D$+$N%H%T%C%/$,(BWeb$B$N$_$K7G<($5$l$k$h$&$K$J$j$^$7$?!#$I$3$K7G<($5$l$k!)(B}
{http://www.debian.or.jp/}
{http://www.debian.org/}
{http://bits.debian.org/}
{C}
{2015/10/22$B$N(BDPN$B$N%a!<%k$+$i!"=>MhN.$l$F$$$?(BDPN$B$N%a!<%k$N9=@.$,:~?7$5$l$^$7$?!#$3$A$i$KH<$$!"%;%-%e%j%F%#$K$D$$$F$N%"%J%&%s%9$H!"?7$7$$(BDD/DM$B$N%"%J%&%s%9$O!"(BDPN$B$N%a!<%k$K$O4^$^$l$:!"(BWeb$B%Z!<%8$K7G:\$5$l$k$N$_$H$J$j$^$7$?!#(B}

\santaku
{2015/11/7$B$K$F!"$H$"$k%$%s%9%?%s%H%a%C%;%s%8%c!<MQ%W%m%H%3%k$N%5!<%S%9$,A4(BDebian Developer$B$G;H$($k$h$&$K$J$C$?$H$N%"%J%&%s%9$,$"$j$^$7$?!#%W%m%H%3%k$NL>A0$O<!$N$I$l!)(B}
{XMPP}
{IRC}
{IP Messanger}
{A}
{XMPP$B$O%*!<%W%s$J%$%s%9%?%s%H%a%C%;%s%8%c!<MQ%W%m%H%3%k$N#1$D!#$J$*!"(BDebian$B4X78<T$NMxMQ$K$"$?$C$F>\$7$/$O!"(Bhttp://rtc.debian.org$B$r;2>H!#(B}

\santaku
{2015/10/30$B$K$F!"(BDebian$B$K$F!"(B2$B2sL\$N8xJg$,9T$o$l$?Lr$^$o$j$O<!$N$I$l!)(B}
{2016 DPN}
{technical commitee}
{2016 Debian JP$B2qD9(B}
{B}
{2$BL>$[$I!"(BDebian Developer$B$NJ}$G(Btechnical comittee$B$G3hLv$G$-$kJ}Jg=8$H$N$3$H$G$9!#8=?&(Btechnical comittee$B$N?M$i$,!"(B2015/12/31$B$GG$4|$,@Z$l$F$7$^$&$H$$$&$3$H$G!"$=$l$^$G$K8uJd<T$r5s$2$kI,MW$,$"$k$H$N$3$H$G$9!#(B}




\section{Debian GNU/kFreeBSD セットアップガイド 2015年版}
\emtext{Debian GNU/kFreeBSD セットアップガイド 2015年版}

\section{今後のイベント}
\emtext{今後のイベント}
\begin{frame}{今後のイベント}
\begin{itemize}
\item 関西エリアDebian勉強会。
\item 12/19(土) 14:00-19:00 第134回東京エリアDebian勉強会\\
  場所はdotsさんにて。
\end{itemize}
\end{frame}

\section{今日の宴会場所}
\emtext{今日の宴会場所}
\begin{frame}{今日の宴会場所}
未定
\end{frame}

\end{document}

;;; Local Variables: ***
;;; outline-regexp: "\\([ 	]*\\\\\\(documentstyle\\|documentclass\\|emtext\\|section\\|begin{frame}\\)\\*?[ 	]*[[{]\\|[]+\\)" ***
;;; End: ***
